% Kalkylark tutorial. Build: make docs   (tectonic)
% Licensed CC BY 4.0. The screens are generated by tools/screens.py.
\documentclass[11pt]{article}
\usepackage{kalkylark}

\title{\sffamily\bfseries\color{ink}Kalkylark\\[4pt]
  \Large\mdseries A tutorial}
\author{\sffamily Set Lonnert and the Kalkylark authors}
\date{\sffamily Version 0.1 · November 2026}

\begin{document}
\maketitle
\thispagestyle{empty}

\begin{center}
\input{budget.tex}
\end{center}

\vspace{1em}
\begin{abstract}
\noindent Kalkylark is a small spreadsheet for text terminals, built the way
spreadsheets were built in 1979. This tutorial goes from an empty sheet to
working models: entering labels and numbers, writing formulas by pointing,
formatting, saving, and the two ways a spreadsheet can recalculate. Every
screen in this document was produced by the program itself.
\end{abstract}

\newpage
\tableofcontents
\newpage

% ======================================================================
\section{Introduction}

In 1979 Dan Bricklin and Bob Frankston released VisiCalc for the Apple~II.
It was the first electronic spreadsheet for a personal computer, and for
many people the reason to buy one. It fitted in 32\,KB together with the
sheet it worked on.

Kalkylark is an independent reconstruction of that program in portable C.
\emph{Kalkylark} is Swedish for \emph{spreadsheet}. It follows VisiCalc's
screen layout, commands, formula language and recalculation rules, and adds
a few things from the decade that followed: natural-order recalculation,
colour themes, and a plain-text file format. It contains no code from the
original.

\subsection{What you need}

A POSIX system (Linux, macOS, a BSD) with a C compiler, and a terminal that
understands VT100 escape sequences, which is all of them. Build and start it
from the source directory:

\begin{lstlisting}
make
./kalkylark -p green
\end{lstlisting}

The \typ{-p} option picks a \emph{phosphor}. Monitors of the time glowed
green, amber or white, and Kalkylark imitates them with \typ{-p green},
\typ{-p amber} and \typ{-p white}. Without \typ{-p} it uses your terminal's
own colours. The screens in this tutorial use green.

% ======================================================================
\section{The screen}

\screen{empty}{An empty sheet.}

The screen has four parts.

\begin{description}[style=nextline, leftmargin=1.5em]
\item[The status line] at the top shows where the \emph{cell cursor} is
  (\cell{A1}), and later what the cell contains. At the right are the
  recalculation order (\textsf{C}, see section~\ref{sec:recalc}) and the
  program's name.
\item[The prompt line] below it holds messages, menus and questions.
\item[The entry line] shows what you are typing.
\item[The sheet] is a grid of \emph{cells}. Columns are named by letters
  \textsf{A} to \textsf{BK} (63 columns), rows by numbers 1 to 254. The
  borders are drawn in reverse video, and so is the cell cursor.
\end{description}

\subsection{Moving around}

The arrow keys move the cell cursor one cell at a time, and the sheet
scrolls when you reach the edge of the screen. \kbd{PgUp} and \kbd{PgDn}
move a screenful, and \kbd{Home} goes back to \cell{A1}.

To jump to a cell, type \typ{>} followed by its coordinate and \enter.
\typ{>BK254}\enter{} takes you to the very last cell.

% ======================================================================
\section{Labels and values}

A cell holds one of two kinds of entry:

\begin{itemize}
\item a \emph{label}: text, such as a heading or the name of a row;
\item a \emph{value}: a number or a formula that computes one.
\end{itemize}

There is no mode to switch. \textbf{The first key you type decides.} A
digit, or one of \typ{+ - . ( @}, starts a value. Anything else starts a
label. Move to \cell{A1} and type \typ{Rent}:

\screen{label}{Typing a label. The prompt line says \textsf{Label}.}

The text appears on the entry line, not yet in the cell. Now press \rgt.
The arrow stores the entry and moves the cursor, so you can fill a row or
column without touching \enter. In \cell{B1} type \typ{9500}:

\screen{value}{A value. The label \textsf{Rent} is now stored in \cell{A1}.}

Press \enter{} to store it, or \rgt/\down{} to store it and move on.
\esc{} throws the entry away, and \kbd{Backspace} edits it.

\begin{note}[Labels that look like numbers]
To store \typ{2026} or \typ{+46 8 123} as text, start with a double quote:
\typ{"2026}. The quote is not stored. It only tells Kalkylark that a label
follows.
\end{note}

Labels are left-aligned, values right-aligned, and a label longer than its
column is cut off at the column edge.

% ======================================================================
\section{Formulas}

Add a second row, \typ{Food} and \typ{4200}, then type \typ{Total} in
\cell{A3}. Now we want \cell{B3} to add up the two amounts.

A formula is a value that refers to other cells. \typ{+B1+B2} would work. The
leading \typ{+} is needed because \typ{B1} on its own starts with a letter
and would become a label. For longer columns, use the function \typ{@SUM}
with a \emph{range}: \typ{@SUM(B1...B2)}.

\subsection{Pointing}

You can type the coordinates, but VisiCalc's way is to \emph{point} at the
cells. Type \typ{@SUM(} in \cell{B3}. Since the entry ends in a \typ{(},
the arrow keys no longer leave the cell. Instead they move a
\emph{pointer}, and its coordinate appears in the formula. Press \up{}
twice:

\screen{pointing}{Pointing. The pointer is on \cell{B1} and the formula
reads \typ{@SUM(B1}.}

Pointing works after an operator, after \typ{(}, and after a comma. Now type
\typ{.}, and Kalkylark writes \typ{...}, VisiCalc's range notation. Point
again, one step \down:

\screen{range}{A range from \cell{B1} to \cell{B2}.}

Type \typ{)} and press \enter:

\screen{sum}{The finished formula. The status line shows the cell's
contents, \textsf{(V)} for value.}

Change \cell{B1} or \cell{B2} and the total follows at once. That instant
recalculation is what made the spreadsheet famous.

\subsection{Left to right: the one rule to remember}

VisiCalc had no operator precedence. Formulas are evaluated \textbf{strictly
from left to right}, and Kalkylark does the same:

\screen{leftright}{\typ{2+3*4} is \textsf{20}: first $2+3=5$, then $5\times4$.}

Most of the time this does what you want, because people tend to write
formulas in the order they think of them. When it doesn't, use parentheses:
\typ{2+(3*4)} is 14. One case deserves a warning. In a loan formula,

\begin{center}
\typ{1-(1+B5)\^{}(-B6)}\quad means\quad $\bigl(1-(1+B_5)\bigr)^{-B_6}$,
\end{center}

which is not what the mathematics says. The power is applied to everything
to its left. Write \typ{1-((1+B5)\^{}(-B6))}. The author of this tutorial
made exactly this mistake in the first draft of the loan example. The result
was \textsf{ERROR}.

\subsection{Operators and functions}

\begin{center}
\begin{tabularx}{\linewidth}{@{}lX@{}}
\toprule
\typ{+ - * /} & add, subtract, multiply, divide (\typ{/0} gives \textsf{ERROR}) \\
\typ{\^{}}      & power: \typ{2\^{}10} is 1024 \\
\typ{-x}       & negation, applied to the operand: \typ{-2\^{}2} is 4 \\
\typ{1.5E3}    & numbers in scientific notation \\
\typ{A1}       & a cell; empty cells and labels count as 0 \\
\typ{A1...C9}  & a range, also written \typ{A1.C9}, \typ{A1..C9} or \typ{A1:C9} \\
\midrule
\typ{@SUM(\dots)}     & sum of values and ranges, e.g. \typ{@SUM(A1...A9,C1,100)} \\
\typ{@MIN}, \typ{@MAX} & smallest, largest \\
\typ{@COUNT}          & number of values (labels and blanks are not counted) \\
\typ{@AVERAGE}        & \typ{@SUM} divided by \typ{@COUNT} \\
\typ{@NPV(r,range)}   & net present value at rate $r$: $\sum v_i/(1+r)^i$ \\
\typ{@LOOKUP(x,range)} & table look-up, see section~\ref{sec:tax} \\
\typ{@ABS @INT @SQRT} & absolute value, integer part, square root \\
\typ{@EXP @LN @LOG10} & exponential and logarithms \\
\typ{@SIN @COS @TAN}  & trigonometry, in radians, also \typ{@ASIN @ACOS @ATAN} \\
\typ{@PI}             & 3.14159\dots \\
\typ{@NA}, \typ{@ERROR} & the special values \textsf{NA} and \textsf{ERROR} \\
\bottomrule
\end{tabularx}
\end{center}

\textsf{NA} (\emph{not available}) marks a number you don't have yet.
Anything computed from it is also \textsf{NA}, so an unfinished model shows
where data is missing instead of showing wrong totals. \textsf{ERROR}
spreads the same way and wins over \textsf{NA}.

If a formula can't be understood (\typ{1+}, say, or an unknown function),
Kalkylark refuses to store it, and the prompt line says what is wrong and
where. Fix it on the entry line and press \enter{} again.

% ======================================================================
\section{Commands}

Everything else is done with \emph{slash commands}. Type \typ{/} and the
prompt line shows the menu:

\screen{menu}{The command menu. Each letter is a command.}

Each command is a letter, sometimes followed by another letter or a value,
and the prompt line always tells you what comes next. \esc{} backs out.
These are available so far:

\begin{center}
\begin{tabularx}{\linewidth}{@{}lX@{}}
\toprule
\typ{/B}\enter   & blank the current cell \\
\typ{/C} \typ{Y} & clear the whole sheet \\
\typ{/F}\emph{f} & format the current cell (section~\ref{sec:formats}) \\
\typ{/GF}\emph{f} & global format \\
\typ{/GC}\emph{n}\enter & column width \\
\typ{/GO}\emph{o} & recalculation order: \typ{C}, \typ{R} or \typ{N} (section~\ref{sec:recalc}) \\
\typ{/SS}       & save the sheet (section~\ref{sec:files}) \\
\typ{/SL}       & load a sheet \\
\typ{/SQY}      & quit \\
\midrule
\typ{>}\emph{cell} & go to a cell \\
\typ{!}         & recalculate now \\
\kbd{Ctrl-L}    & redraw the screen \\
\bottomrule
\end{tabularx}
\end{center}

% ======================================================================
\section{Formats and column width}\label{sec:formats}

A \emph{format} changes how a cell's value is shown, not the value itself.
Here the same number, 12.345, appears in each format, \textsf{G} to
\textsf{*} across \cell{A1}--\cell{F1}:

\screen{formats}{One value, six formats, in 7-character columns.}

\begin{center}
\begin{tabularx}{\linewidth}{@{}clX@{}}
\toprule
\typ{G} & general & as many digits as fit in the column \\
\typ{I} & integer & rounded to a whole number \\
\typ{\$} & dollars & two decimals, for money in any currency \\
\typ{L} & left    & left-aligned \\
\typ{R} & right   & right-aligned; also works for labels \\
\typ{*} & graph   & one star per unit: a bar chart made of cells \\
\typ{D} & default & remove the cell's format, use the global one \\
\bottomrule
\end{tabularx}
\end{center}

\typ{/F}\emph{f} formats the current cell. A cell keeps its format when you
type new contents into it. \typ{/GF}\emph{f} sets the \emph{global} format,
used by every cell that has no format of its own. When a number does not
fit in its column, the cell fills with \typ{>>>>}. Widen the columns, or
choose a shorter format.

\typ{/GC} sets the width of all columns, from 3 to 40 characters. The
current width is suggested: press \enter{} to keep it, or type a new
number. As in VisiCalc, every column has the same width.

% ======================================================================
\section{Saving and loading}\label{sec:files}

\typ{/SS} saves the sheet and asks for a file name. \typ{/SL} loads one,
replacing the current sheet. You can also name the file when you start:

\begin{lstlisting}
./kalkylark budget.kalk
\end{lstlisting}

If the file exists it is loaded. Otherwise it is the name the first
\typ{/SS} will use. Kalkylark remembers the name, so later saves only need
\enter. It asks before overwriting a file, and before quitting or loading
over unsaved changes.

The file is plain text, so you can read it, compare versions, or generate
it from a script:

\begin{lstlisting}
# kalkylark 1
/GO N
/GC 11
/GF G
>A3:"Rent
>B3:/F$:9500
>B8:/F$:@SUM(B3...B6)
\end{lstlisting}

The \typ{/G} lines are the global settings. Each \typ{>} line is one cell:
its coordinate, an optional format, and the contents exactly as you would
type them. Kalkylark saves to a temporary file first and then renames it, so
a failed save never damages the old file.

% ======================================================================
\section{Recalculation}\label{sec:recalc}

When you change a cell, every formula is computed again. The \emph{order}
matters more than you might think.

\subsection{VisiCalc's way: one sweep}

VisiCalc went through the sheet once, column by column (or row by row), and
computed each formula from whatever its neighbours held at that moment. Try
this: in \cell{A1} type \typ{+A2*2}, in \cell{A2} \typ{+A3+1}, and in
\cell{A3} the number \typ{5}.

\screen{stale}{Column order: \cell{A1} shows 2, but $(5+1)\times2=12$.}

\cell{A1} was computed first, while \cell{A2} still held its old value. The
formulas refer \emph{forward}, to cells that come later in the sweep. Press
\typ{!} to recalculate again and the numbers settle. VisiCalc users learned
to press \typ{!} until nothing changed, and to arrange their sheets so that
formulas refer up and to the left.

\subsection{Natural order}

Lotus 1-2-3 (1983) introduced \emph{natural order}, which Kalkylark offers
as \typ{/GON}. Each formula is computed after all the formulas it reads,
wherever they are on the sheet:

\screen{natural}{After \typ{/GON}: one recalculation, right answer. The
status line shows \textsf{N}.}

\typ{/GOC} and \typ{/GOR} go back to VisiCalc's column and row order. The
order is saved with the sheet. Natural order is the practical choice, while
column order is the authentic one.

\subsection{Circular references}

Sometimes formulas depend on each other in a circle. Here a fee is 10\% of
a total that includes the fee (\typ{examples/iterate.kalk}):

\screen{circ}{\cell{B2} is \typ{+B3*.1} and \cell{B3} is \typ{+B1+B2}.
The status line shows \textsf{CIRC}.}

No order can satisfy both formulas at once, so in natural order Kalkylark
warns with \textsf{CIRC}. It still computes the cells in the circle, one
sweep per recalculation. Press \typ{!} a few times and the numbers
converge:

\screen{circ-settled}{After five presses of \typ{!}: close to 111.11 and
11.11.}

Usually \textsf{CIRC} means a mistake, such as a total that includes
itself. Then the fix is to correct the formula, and the warning goes away.

% ======================================================================
\section{Worked examples}

The \typ{examples/} directory contains the sheets used below. Open one by
naming it when you start, for example:

\begin{lstlisting}
./kalkylark examples/budget.kalk
\end{lstlisting}

\subsection{A household budget}

\typ{budget.kalk} is the sheet on the title page: labels in column
\textsf{A}, amounts in \textsf{B} formatted with \typ{/F\$}, a total
\typ{@SUM(B3...B6)}, a daily figure \typ{+B8/30}, and the largest item
\typ{@MAX(B3...B6)}. Try changing the rent and watch the total, the daily
figure and the maximum change together.

\subsection{A loan and an investment}

\typ{loan.kalk} computes the monthly payment on a loan of 200\,000 at 4\%
a year over 25 years:

\screen{loan}{The payment formula, with its extra parentheses.}

\begin{center}
\begin{tabular}{@{}lll@{}}
\toprule
\cell{B5} & \typ{+B2/12} & monthly rate \\
\cell{B6} & \typ{+B3*12} & number of payments \\
\cell{B7} & \typ{+B1*B5/(1-((1+B5)\^{}(-B6)))} & payment: 1055.67 \\
\cell{B8} & \typ{+B7*B6} & total paid \\
\bottomrule
\end{tabular}
\end{center}

Further down, an investment of 10\,000 that returns 3\,000 a year for five
years is valued with \typ{@NPV(B17,B11...B15)+B10}. At a discount rate of
8\% it is worth 1978.13 today, so it is worth doing.

\subsection{A tax table}\label{sec:tax}

\typ{@LOOKUP(x,range)} searches a column (or row) of thresholds in
ascending order, finds the last one that is not greater than $x$, and
returns the value \emph{beside} it: to the right for a column, below for a
row. In \typ{tax.kalk} the thresholds are in \cell{A2}--\cell{A5} and the
rates in \textsf{B}:

\screen{tax}{\typ{@LOOKUP(E1,A2...A5)}: an income of 65\,000 falls in the
50\,000 bracket, rate 0.2.}

If $x$ is smaller than the first threshold, the result is \textsf{NA}.

% ======================================================================
\section{Where to go next}

\begin{itemize}
\item \typ{man ./kalkylark.1} is the complete reference, in the terminal.
\item The \emph{history} document describes how spreadsheets learned to
  recalculate, from VisiCalc's sweep to dependency graphs, and the legal
  background of a reconstruction like this one.
\item The \emph{internals} document walks through the source code: the
  terminal layer, the formula parser and the recalculation engine.
\end{itemize}

Kalkylark is free software under the BSD 2-clause licence, and this
tutorial is licensed CC BY 4.0. VisiCalc is a trademark of its respective
owner. Kalkylark is not affiliated with Lotus, IBM, Software Arts or
VisiCorp.

\end{document}
